\documentclass[11pt,a4paper]{article}
\usepackage[margin=2.2cm]{geometry}
\usepackage{amsmath,amssymb}
\usepackage{graphicx}
\usepackage{booktabs}
\usepackage{xcolor}
\usepackage[colorlinks=true,linkcolor=blue!50!black,urlcolor=blue!50!black]{hyperref}
\usepackage{enumitem}
\usepackage{fontspec}
\setmonofont{Menlo}[Scale=0.82]
\setlist{itemsep=2pt,topsep=4pt}

\newcommand{\code}[1]{\texttt{#1}}

\title{Immersed interface method in StagFDTools\\[2pt]
\large Notes on steps 1--2: interface geometry and interface velocity}
\author{StagFDTools, branch \texttt{iim-discontinuous-viscosity}}
\date{\today}

\begin{document}
\maketitle

\noindent This note documents the first steps of the implementation plan
(\texttt{docs/IIM\_plan.md}) for porting the immersed interface method (IIM) of
W.~Landry (\url{https://gitlab.com/wlandry/discontinuous_stokes}) to StagFDTools.
In that method, the Stokes equations are discretised on the standard MAC grid in terms
of $z=\eta\mathbf v$ and $p$. Stencils that cross the viscosity interface $\Gamma$ receive
correction terms built from the jumps $[z]$, $[\nabla z]$ and $[p]$.
\begin{itemize}
  \item \textbf{Step~1} (Sections~\ref{sec:s1impl}--\ref{sec:s1res}) provides the geometry
    these corrections need: where $\Gamma$ crosses the grid, the normal there, and which
    equations are affected.
  \item \textbf{Step~2} (Section~\ref{sec:s2}) provides the interface velocity
    $\mathbf v_\Gamma$ and its tangential derivative $\partial_t\mathbf v_\Gamma$, on which
    all the jumps depend.
\end{itemize}
The code is in \code{src/StokesIIM/}, exported by the submodule \code{StagFDTools.StokesIIM}.

%---------------------------------------------------------------
\section{Discretisation}

Figure~\ref{fig:mesh} summarises the discretisation. There are no unknowns beyond the
standard ones: $z_x$ on the $x$-faces, $z_y$ on the $y$-faces and $p$ at the centroids.

An equation is \emph{irregular} if its stencil has members on both sides of $\Gamma$. For each
such equation, the jumps are evaluated where $\Gamma$ crosses the segment between stencil
members, and they depend on two interface quantities:
\begin{itemize}
  \item the interface velocity $\mathbf v_\Gamma$;
  \item its tangential derivative $\partial_t\mathbf v_\Gamma$.
\end{itemize}
These are not solved for. They are eliminated by one-sided linear fits on each side, combined
with a viscosity-weighted average:
\[
\mathbf v_\Gamma=\frac{\eta_-\mathbf v_-+\eta_+\mathbf v_+}{\eta_-+\eta_+}.
\]

\begin{figure}[h]
  \centering
  \includegraphics[width=\textwidth]{figures/iim_mesh.png}
  \caption{IIM discretisation on the MAC grid (script \code{docs/figures/iim\_mesh.jl}).
  \textbf{(a)}~Unknowns coloured by phase (orange: inclusion $\eta_-$, blue: matrix $\eta_+$).
  Outlined nodes are irregular equations. Gold diamonds are the crossings of $\Gamma$ with the
  half-grid lines, where jumps are evaluated; they are not unknowns.
  \textbf{(b)}~One irregular $x$-momentum equation. The star and hexagon are the interface
  points it uses for $\partial_{xx}z_x,\partial_{yy}z_x$ and for the cut corner of
  $\partial_{xy}z_y$. Crosses mark the fit centres $x_\Gamma\pm L\mathbf n$. Triangles show the
  3-node one-sided fits (solid: $z_x$, dashed: $z_y$) that give $\mathbf v_\Gamma$ and
  $\partial_t\mathbf v_\Gamma$.}
  \label{fig:mesh}
\end{figure}

%---------------------------------------------------------------
\section{Step 1: implementation}\label{sec:s1impl}

\paragraph{Convention.} $\Gamma=\{\varphi=0\}$, with $\varphi<0$ inside the inclusion and
$\varphi>0$ in the matrix. The function \code{side($\varphi$)} (line~14) returns $-1$ or $+1$.
It never returns 0, so every node belongs to exactly one phase.

\begin{enumerate}[leftmargin=*]
\item \textbf{Sampling the level set} -- \code{LevelSet(φfun, X)} (line~22).
  It evaluates $\varphi$ on the extended grids of \code{GenerateGrid}, ghost nodes included,
  and returns the NamedTuple \code{(x, y, c, v)} for the $V_x$, $V_y$, centroid and vertex
  arrays, following the StagFDTools convention.
  \begin{itemize}
    \item Nodes with $|\varphi|<10^{-10}\Delta$ lie on $\Gamma$; they are moved to the $+$ side.
    \item Without this, a node exactly on $\Gamma$ (e.g.\ the vertex $(0,2)$ for $r_c=2$,
      $\Delta=0.5$) gets $\varphi=-7\cdot10^{-16}$ by round-off. That creates an inside sliver
      of zero width, which stencils cross twice.
  \end{itemize}

\item \textbf{Fine lattice} -- \code{FineLattice(φ, X, Δ)} (line~38).
  The four staggered arrays are merged into one array $\Phi_{IJ}$ with spacing $\Delta/2$:
  \begin{center}
  \begin{tabular}{@{}ll@{}}
    \toprule
    $(I,J)$ parity & node \\
    \midrule
    odd/odd & vertex \\
    even/even & centroid \\
    odd/even & $V_x$ \\
    even/odd & $V_y$ \\
    \bottomrule
  \end{tabular}
  \end{center}
  Every interface point used by an IIM stencil lies on an edge of this lattice:
  \begin{itemize}
    \item $V_x$--$V_x$ and $V_x$--$p$ segments;
    \item the box edges of the mixed derivative;
    \item the centroid--face segments of continuity.
  \end{itemize}
  The crossings are therefore computed once, on the lattice edges, rather than separately for
  each equation.

\item \textbf{Crossings} -- \code{InterfaceCrossings(fl)} (line~109).
  On every lattice edge where \code{side} changes:
  \begin{itemize}
    \item build the cubic Lagrange polynomial through the 4 samples $k-1,\dots,k+2$ along the
      grid line (functions \code{crossing\_fraction}, line~79, and \code{lagrange}, line~61);
    \item near the ends of the lattice, where only 3 samples exist, use the quadratic instead;
    \item locate the root on $[0,1]$ by bisection. This always succeeds, because the
      polynomial takes the endpoint values, which have opposite signs.
  \end{itemize}
  The output is:
  \begin{itemize}
    \item \code{tx}, \code{ty}: the fraction along each horizontal or vertical edge
      (\code{NaN} if the edge is not crossed);
    \item \code{points}: a list of $(x_\Gamma,\mathbf n,I,J,\text{dir})$, one per crossing.
  \end{itemize}

\item \textbf{Normals} -- function \code{∇Φ} (line~93).
  \begin{itemize}
    \item At both ends of the edge, $\nabla\varphi$ is approximated by central differences
      (lattice spacing $\Delta/2$, so a stencil width of $\Delta$).
    \item These are interpolated linearly to $x_\Gamma$ and normalised.
    \item $\mathbf n$ points towards $\varphi>0$, i.e.\ out of the inclusion.
  \end{itemize}

\item \textbf{Irregular equations} -- \code{IrregularNodes(φ, nc)} (line~145).
  An equation is flagged when any member of its standard stencil lies on the other side of
  $\Gamma$. The indexing matches \code{Stokes.jl}:
  \begin{center}
  \begin{tabular}{@{}ll@{}}
    \toprule
    equation & stencil members checked \\
    \midrule
    $V_x[i,j]$ & $V_x[i\pm1,j]$, $V_x[i,j\pm1]$, $P[i-1{:}i,\,j-1]$, $V_y[i{:}i+1,\,j-1{:}j]$ \\
    $V_y[i,j]$ & $V_y[i\pm1,j]$, $V_y[i,j\pm1]$, $P[i-1,\,j-1{:}j]$, $V_x[i-1{:}i,\,j{:}j+1]$ \\
    $P[i,j]$ & $V_x[i{:}i+1,\,j+1]$, $V_y[i+1,\,j{:}j+1]$ \\
    \bottomrule
  \end{tabular}
  \end{center}

\item \textbf{Thin features} -- \code{MultipleCrossings(Γ)} (line~177).
  Counts the lattice segments of length $\Delta$ that $\Gamma$ crosses twice. These are
  features thinner than about $\Delta$, which the first-order IIM does not handle (Landry's
  ``anticorners''). The count is $0$ at every resolution of the inclusion test.
\end{enumerate}

\noindent Unit tests are in \code{test/test\_iim\_geometry.jl}, which is included in
\code{runtests.jl}. They check that:
\begin{itemize}
  \item the fine lattice reproduces $\varphi$ at every node;
  \item crossings and normals are exact (to $10^{-12}$) for a straight inclined interface;
  \item the irregular columns for a vertical interface match a hand check;
  \item nodes on $\Gamma$ are moved to the $+$ side.
\end{itemize}
The example is \code{examples/\_Stokes\_IIM/IIM\_Step1\_Geometry.jl}.

%---------------------------------------------------------------
\section{Step 1: results for the viscous inclusion geometry}\label{sec:s1res}

The test geometry is that of Schmid \& Podladchikov (2003): a circular inclusion of radius
$r_c=2$ centred in a square domain of side $L=10$. The grid spacing is
$h=\Delta x=\Delta y=L/n_c$, and two level-set functions with the same zero contour are used:
\[
\varphi_d=r-r_c\quad\text{(signed distance)},\qquad
\varphi_t=\tanh(r-r_c)\quad\text{(not a distance function)}.
\]

\begin{figure}[h]
  \centering
  \includegraphics[width=\textwidth]{../examples/_Stokes_IIM/IIM_Step1_Geometry.png}
  \caption{Geometry for the inclusion test.
  \textbf{Left:} $n_c=20$, $\varphi=\varphi_d$. Circles are $P$, right-pointing triangles $V_x$,
  upward triangles $V_y$; colours give the phase. Outlined nodes are the irregular equations.
  Gold diamonds are the 60 crossings of $\Gamma$ with the fine lattice.
  \textbf{Middle:} zoom on the first quadrant with the computed normals $\mathbf n$.
  \textbf{Right:} maximum error of crossing positions and normals against $h$, for $\varphi_d$
  and $\varphi_t$, with $O(h^2)$ and $O(h^4)$ reference slopes.}
  \label{fig:step1}
\end{figure}

\paragraph{Left and middle panels.}
\begin{itemize}
  \item The irregular equations form a band one stencil wide on each side of $\Gamma$: 48
    $V_x$, 48 $V_y$ and 24 $P$ equations at $n_c=20$.
  \item Only $V_x$ and $V_y$ nodes up to one cell away from $\Gamma$ are flagged, because the
    momentum stencils reach $\pm h$. Only cells cut by $\Gamma$ between their centroid and a
    face are flagged for continuity.
  \item The normals are perpendicular to the circle and point outward.
  \item Some crossings come in close pairs near grid nodes. These are crossings of the
    horizontal and the vertical lattice line through the same region, which is expected.
\end{itemize}

\subsection*{Discussion of the convergence (right panel)}

The errors are measured over all crossings inside the physical domain:
\[
e_x=\max_\Gamma\bigl|\,|x_\Gamma|-r_c\bigr|,\qquad
e_n=\max_\Gamma\bigl|\mathbf n-x_\Gamma/|x_\Gamma|\bigr|.
\]
Table~\ref{tab:conv} lists the values and the local orders $\log_2(e_{2h}/e_h)$.

\begin{table}[h]
\centering
\small
\begin{tabular}{@{}rrrr rr rr rr rr@{}}
\toprule
& & & & \multicolumn{4}{c}{$\varphi_d=r-r_c$} & \multicolumn{4}{c}{$\varphi_t=\tanh(r-r_c)$}\\
\cmidrule(lr){5-8}\cmidrule(l){9-12}
$n_c$ & $h$ & points & irr.\ $V_x/V_y/P$ & $e_x$ & ord. & $e_n$ & ord. & $e_x$ & ord. & $e_n$ & ord.\\
\midrule
 10 & 1      &   28 & 24/24/12    & 3.9e-4  & --   & 9.3e-3 & --   & 1.5e-3 & --   & 1.5e-2 & --   \\
 20 & 0.5    &   60 & 48/48/24    & 1.9e-5  & 4.39 & 2.4e-3 & 1.94 & 7.6e-5 & 4.28 & 4.4e-3 & 1.74 \\
 40 & 0.25   &  124 & 88/88/44    & 1.7e-6  & 3.44 & 7.9e-4 & 1.63 & 6.8e-6 & 3.48 & 1.2e-3 & 1.93 \\
 80 & 0.125  &  252 & 184/184/92  & 1.0e-7  & 4.08 & 2.1e-4 & 1.89 & 3.8e-7 & 4.18 & 2.8e-4 & 2.06 \\
160 & 0.0625 &  508 & 360/360/180 & 6.6e-9  & 3.94 & 4.9e-5 & 2.13 & 2.7e-8 & 3.79 & 7.3e-5 & 1.91 \\
320 & 0.03125& 1020 & 728/728/364 & 4.3e-10 & 3.95 & 1.3e-5 & 1.90 & 1.7e-9 & 4.00 & 1.8e-5 & 2.00 \\
\bottomrule
\end{tabular}
\caption{Convergence of interface positions and normals for the inclusion ($r_c=2$, $L=10$).
\emph{points}: number of crossings; \emph{irr.}: number of irregular equations.
No multiple crossings occur at any resolution.}
\label{tab:conv}
\end{table}

\paragraph{Positions converge at fourth order.}
\begin{itemize}
  \item Along a grid line, the crossing is the root of the cubic interpolant through samples
    spaced $\Delta/2$. The interpolation error is
    \[
    \frac{\varphi^{(4)}}{4!}\prod_k(s-s_k)=O(h^4),
    \]
    and bisection adds nothing (60 halvings).
  \item The error that matters is the distance of $x_\Gamma$ from $\Gamma$. Shifting the root
    along the line by $\delta s=\delta\varphi/|\partial_s\varphi|$ moves the point off
    $\Gamma$ by
    \[
    \delta s\,|\mathbf n\cdot\mathbf e_s|=\delta\varphi/|\nabla\varphi|.
    \]
  \item The measured error $e_x$, a distance normal to $\Gamma$, is therefore \emph{not}
    amplified at grazing intersections,
    where $\Gamma$ is nearly parallel to the grid line. Only the tangential position of the
    point is less accurate there, and the point still lies on $\Gamma$ to $O(h^4)$.
  \item The local orders oscillate between 3.4 and 4.4 at coarse resolution. As $h$ changes,
    the crossings move relative to the samples, and the factor $\prod_k(s-s_k)$ depends on
    where in the stencil the root lies. The maximum is taken over a different set of points at
    each resolution. The mean slope is 4.
\end{itemize}

\paragraph{Normals converge at second order.}
Two approximations each contribute an $O(h^2)$ error:
\begin{itemize}
  \item the central difference over a width $\Delta$, with truncation error
    $(\Delta^2/6)\,\partial^3\varphi$;
  \item the linear interpolation of $\nabla\varphi$ along an edge of length $\Delta/2$.
\end{itemize}
Even for the exact distance function, $\partial^3\varphi\neq0$ because of the curvature
$\kappa=1/r_c$. Hence $e_n\gg e_x$ in absolute terms ($10^{-5}$ against $10^{-10}$ at
$n_c=320$). The order is nevertheless clean (1.6--2.1).

\paragraph{A distance function is not required.}
$\varphi_t$ gives the same orders as $\varphi_d$, with constants larger by a factor of about
1.5--4. Both functions have $|\nabla\varphi|=1$ on $\Gamma$, so the larger constants come from
the larger higher derivatives of $\tanh$. This matters for later steps: a level set reconstructed from
markers will not be an exact distance function, and no redistancing is needed for the geometry
to converge.

\paragraph{Consequences for the IIM.}
\begin{itemize}
  \item In the first-order IIM, the corrections have the form
    $C\sim([z]+[\nabla z]\cdot\boldsymbol\delta)/h^2$ and $C_p=\pm[p]/h$, where
    $\boldsymbol\delta$ is the distance from $x_\Gamma$ to a grid node.
  \item \emph{Position error.} An error $e_x=O(h^4)$ in $\boldsymbol\delta$ perturbs $C$ by
    $O(h^2)$, which is negligible.
  \item \emph{Normal error.} An error $e_n=O(h^2)$ rotates the $(n,t)$ frame in which the jumps
    are written. This perturbs $[\nabla z]$ and $[p]=-2[\eta]\,\partial_t v_t$ by
    $O(h^2)[\eta]$, which contributes $O(h)$ to the truncation error of the irregular
    equations.
  \item The one-sided linear fits give $\mathbf v_\Gamma$ to $O(h^2)$ and
    $\partial_t\mathbf v_\Gamma$ to $O(h)$. Once divided by $h^2$ (and multiplied by
    $\boldsymbol\delta\sim h$ for the derivative term), this leaves an $O(1)$ truncation error at
    the irregular equations. That is the reason the scheme is only first order overall.
  \item The geometry contributes at most $O(h)$, one order below the interpolation error, so it
    is not expected to limit the accuracy of the first-order scheme. The second-order variant
    will also need the curvature and the second derivatives of $\varphi$. These can be obtained
    by central differences on the same lattice.
  \item The irregular equations scale linearly with $n_c$ (about $2.3\,n_c$ for each
    velocity component), so the extra cost of the corrections is $O(n_c)$ against $O(n_c^2)$
    unknowns.
\end{itemize}

%---------------------------------------------------------------
\section{Step 2: interface velocity and tangential derivative}\label{sec:s2}

\paragraph{Method.} At each crossing $x_\Gamma$ (normal $\mathbf n$, tangent
$\mathbf t=(-n_y,n_x)$), and for each velocity component:
\begin{itemize}
  \item On each side $s=\pm1$, the 3 nodes of that component's grid closest to the fit centre
    $x_\Gamma+sL\mathbf n$ define a linear interpolant of $z$. Here
    $L=1/\max(|n_x|/\Delta x,|n_y|/\Delta y)$ is one cell along $\mathbf n$, as in Landry's
    \code{compute\_dv\_dtt.cxx}.
  \item The interpolant's value and gradient at $x_\Gamma$ give $z_s(x_\Gamma)$ and $\partial_t z_s$.
  \item $\eta$ is constant on each side and $v_s=z_s/\eta_s$, so Landry's viscosity-weighted
    average reduces to
    \[
    \mathbf v_\Gamma=\frac{z_-(x_\Gamma)+z_+(x_\Gamma)}{\eta_-+\eta_+},\qquad
    \partial_t\mathbf v_\Gamma=\frac{\partial_t z_-+\partial_t z_+}{\eta_-+\eta_+}.
    \]
  \item This is exact when $\mathbf v$ is linear on each side and continuous across $\Gamma$
    (then $z_\pm=\eta_\pm\mathbf v$).
  \item Both quantities are linear in the 6 nodal values of $z$. They are therefore stored as
    weights, which the assembly of step~4 will turn into matrix coefficients.
  \item \emph{Deviation from Landry:} he takes $\partial_t z_\pm$ from separate nearest
    one-sided differences. Here it comes from the gradient of the same fit. Both are $O(h)$.
\end{itemize}

\paragraph{Implementation} (\code{src/StokesIIM/Trace.jl}).
\begin{enumerate}[leftmargin=*]
\item \code{OneSidedFit} (line~20).
  \begin{itemize}
    \item Collect the candidate nodes within $\pm4$ cells of the fit centre that lie on side
      $s$, and sort them by distance.
    \item Take the 2 closest, then the next one that is not collinear with them
      ($|\det[1,\,\delta x,\,\delta y]|>1/2$ in cell units). Unlike Landry's row/column test,
      this also rejects diagonal alignments.
    \item Return the barycentric weights $\lambda_k$ of $x_\Gamma$ and the gradient weights
      $\mathbf g_k$, both from the inverse of the $3\times3$ matrix (\code{StaticArrays}).
  \end{itemize}
\item \code{InterfaceTrace} (line~53).
  \begin{itemize}
    \item Run the two one-sided fits.
    \item Return the 6 node indices and the weights $v_k=\lambda_k/(\eta_-+\eta_+)$ and
      $t_k=\mathbf g_k\cdot\mathbf t/(\eta_-+\eta_+)$, so that
      $v_\Gamma=\sum_k v_kz_k$ and $\partial_tv_\Gamma=\sum_k t_kz_k$.
  \end{itemize}
\item \code{InterfaceTraces} (line~72). Applies \code{InterfaceTrace} to both components at
  every crossing inside the domain. Candidates are the $V_x$ and $V_y$ nodes of the domain,
  boundary nodes included (the ranges of \code{Ranges}).
\item \code{apply\_trace} (line~85). Evaluates the weights on an array $z$.
\end{enumerate}

\noindent Tests are in \code{test/test\_iim\_trace.jl}. They check that:
\begin{itemize}
  \item the weights are exact (to $10^{-10}$) for a velocity field with a kink across an
    inclined $\Gamma$;
  \item $\sum_k v_k\eta_k=1$ and $\sum_k t_k\eta_k=0$;
  \item each fit uses only nodes of its own side.
\end{itemize}
The example is \code{examples/\_Stokes\_IIM/IIM\_Step2\_Trace.jl}. It samples
$z=\eta\mathbf v$ from the exact solution and compares the reconstructed values with the exact
ones:
\begin{itemize}
  \item $\mathbf v(x_\Gamma)$ for the value;
  \item $\mathbf L_{inc}\mathbf t$ for the derivative, where $\mathbf L_{inc}$ is the uniform
    velocity gradient of the inclusion ($\partial_t\mathbf v$ is continuous across $\Gamma$).
\end{itemize}

\begin{figure}[h]
  \centering
  \includegraphics[width=\textwidth]{../examples/_Stokes_IIM/IIM_Step2_Trace.png}
  \caption{Interface velocity for the inclusion test ($r_c=2$, $L=10$, $\eta_{mat}=1$).
  \textbf{Left:} the fits selected by the code at three crossings ($n_c=20$,
  $\eta_{inc}/\eta_{mat}=10^3$). Stars are the crossings $x_\Gamma$ and crosses the fit centres
  $x_\Gamma\pm L\mathbf n$. Triangles are the 3-node fits on the inclusion side (orange) and the
  matrix side (blue), solid for $V_x$ and dashed for $V_y$.
  \textbf{Middle, right:} maximum error of $v_\Gamma$ and $\partial_t v_\Gamma$ over all
  crossings and both components against $h$, for three viscosity contrasts.}
  \label{fig:step2}
\end{figure}

\begin{table}[h]
\centering
\small
\begin{tabular}{@{}rr rr rr rr@{}}
\toprule
& & \multicolumn{4}{c}{$\eta_{inc}/\eta_{mat}=10^{-3}$} & \multicolumn{2}{c}{$10^{6}$}\\
\cmidrule(lr){3-6}\cmidrule(l){7-8}
$n_c$ & $h$ & $e_v$ & ord. & $e_t$ & ord. & $e_v$ & ord.\\
\midrule
  20 & 0.5     & 5.4e-1 & --   & 8.2e-1 & --   & 5.6e-5  & --   \\
  40 & 0.25    & 1.9e-1 & 1.52 & 6.2e-1 & 0.41 & 5.1e-6  & 3.47 \\
  80 & 0.125   & 6.9e-2 & 1.43 & 4.0e-1 & 0.64 & 2.6e-7  & 4.28 \\
 160 & 0.0625  & 2.2e-2 & 1.62 & 2.6e-1 & 0.62 & 2.1e-8  & 3.66 \\
 320 & 0.03125 & 6.2e-3 & 1.85 & 1.6e-1 & 0.69 & 6.2e-9  & 1.73 \\
 640 & 0.0156  & 1.7e-3 & 1.86 & 9.1e-2 & 0.81 & 1.7e-9  & 1.90 \\
1280 & 0.0078  & 4.5e-4 & 1.93 & 4.7e-2 & 0.97 & 4.5e-10 & 1.90 \\
\bottomrule
\end{tabular}
\caption{Errors of $v_\Gamma$ ($e_v$) and $\partial_tv_\Gamma$ ($e_t$). For the contrast $10^3$,
all values equal those for $10^{-3}$ times $10^{-3}$, with identical orders. For $10^6$, $e_t$
equals the $10^{-3}$ values times $10^{-6}$.}
\label{tab:step2}
\end{table}

\paragraph{Discussion.}
\begin{itemize}
  \item \emph{Orders.} $v_\Gamma$ converges at $O(h^2)$ and $\partial_t v_\Gamma$ at $O(h)$,
    as expected for linear fits. This holds only asymptotically: the derivative reaches order~1
    only around $n_c=1280$, about 250 cells per radius. Below that, the observed order is
    0.4--0.8, and the error behaves as $e_t\approx h\,|\nabla^2\mathbf v|$ on the matrix side,
    where $|\nabla^2\mathbf v|\approx12/r_c$ next to $\Gamma$.
  \item \emph{Dependence on contrast.} Inside the inclusion the exact velocity is linear, so the
    inclusion-side fit is exact. The whole error therefore comes from the matrix-side fit,
    multiplied by its weight $\eta_{mat}/(\eta_{inc}+\eta_{mat})$. This explains the factors of
    $10^{-3}$ and $10^{-6}$ between the curves.
  \item \emph{What the IIM uses.} The quantity entering the corrections is the jump
    $[z]=\Delta\eta\,v_\Gamma$, with error $|\Delta\eta|\,e_v$. This error is the same for the
    contrasts $10^{-3}$ and $10^3$ (0.54 at $n_c=20$). For $10^6$ it is the same from
    $n_c\approx320$ on ($6.2\cdot10^{-3}$), once the artifact discussed below has decayed. The
    interface treatment therefore does not degrade at high contrast.
  \item \emph{Coarse $10^6$ case.} Here $e_v$ is larger than the contrast scaling predicts and
    decays at about $O(h^4)$. This is the position error of step~1, not an interpolation error:
    $x_\Gamma$ lies $e_x=O(h^4)$ off $\Gamma$, and the exact velocity has a kink there, so the
    reference value changes by $[\partial_n v]\,e_x$. It becomes visible only once the fit
    error has been multiplied by $10^{-6}$.
  \item \emph{Consequence.} The tangential derivative is the least accurate ingredient. It sets
    $[\nabla z]$ and $[p]=-2[\eta]\partial_tv_t$, and therefore the pressure jump. If step~5
    shows that it limits the pressure accuracy, the remedy is a 6-point quadratic fit (Landry's
    second-order version) or a least-squares fit over more nodes.
\end{itemize}

%---------------------------------------------------------------
\section{Next steps}

Step 3 adds the jump conditions ($[z]$, $[\nabla z]$, $[p]$ from $\mathbf v_\Gamma$ and
$\partial_t\mathbf v_\Gamma$). Step 4 adds the corrected momentum and continuity kernels and a
direct solve. The goal is the viscous inclusion test: unlike the standard averaging approach,
the IIM pressure should converge in $L_\infty$.

\end{document}
